% problem_id: S001-lemma-lqf-separated-shriek-composition
% source_id: S001 (Stacks Project)
% source_file: more-etale.tex
% statement_locator: more-etale.tex lines 1401--1418
% proof_locator: more-etale.tex lines 1420--1608
% env: lemma
% id_note: label 在本来源内唯一
% pairing: tight (verified)
% status: complete (statement + author proof verbatim + reason + locators)
%
\begin{lemma}
\label{lemma-lqf-separated-shriek-composition}
Let $f' : X \to Y'$ and $g : Y' \to Y$ be composable morphisms of schemes
with $f'$ and $f = g \circ f'$ locally quasi-finite and $g$ separated and
locally of finite type. Then there is a canonical isomorphism of functors
$g_! \circ f'_! = f_!$. This isomorphism is compatible with
\begin{enumerate}
\item[(a)] covariance with respect to open embeddings as in
Remarks \ref{remark-covariance-f-shriek-separated} and
\ref{remark-covariance-lqf-f-shriek},
\item[(b)] the base change isomorphisms of
Lemmas \ref{lemma-lqf-base-change-f-shriek}
and \ref{lemma-base-change-f-shriek-separated}, and
\item[(c)] equal to the isomorphism of Lemma \ref{lemma-f-shriek-composition}
via the identifications of Lemma \ref{lemma-finite-support-f-shriek-separated}
in case $f'$ is separated.
\end{enumerate}
\end{lemma}

\begin{proof}
Let $\mathcal{F}$ be an abelian sheaf on $X_\etale$. With conventions as in
Remark \ref{remark-construct-map-presheaves-downstairs} we will explicitly
construct a map
$$
c : f_{p!}\mathcal{F} \longrightarrow g_*f'_{p!}\mathcal{F}
$$
of abelian presheaves on $Y_\etale$. By the discussion in
Remark \ref{remark-construct-map-presheaves-downstairs}
this will determine a canonical map
$c^\# : f_!\mathcal{F} \to g_*f'_!\mathcal{F}$.
We will show that $c^\#$ has image contained in the subsheaf
$g_!f'_!\mathcal{F}$, thereby obtaining a map
$c' : f_!\mathcal{F} \to g_!f'_!\mathcal{F}$. Next, we will prove
(a), (b), and (c) that. Finally, part (b)
will allow us to show that $c'$ is an isomorphism.

\medskip\noindent
Construction of the map $c$. Let $V \in Y_\etale$ and
let $s = \sum (Z_i, s_i)$ be a sum as in (\ref{equation-formal-sum})
defining an element of $f_{p!}\mathcal{F}(V)$.
Recall that $Z_i \subset X_V = X \times_Y V$
is a locally closed subscheme finite over $V$.
Setting $V' = Y' \times_Y V$ we get $X_{V'} = X \times_{Y'} V' = X_V$.
Hence $Z_i \subset X_{V'}$ is locally closed and
$Z_i$ is finite over $V'$ because $g$ is separated
(Morphisms, Lemma \ref{morphisms-lemma-finite-permanence}).
Hence we may set $c(s) = \sum (Z_i, s_i)$ but now viewed
as an element of $f'_{p!}\mathcal{F}(V') = (g_*f'_{p!}\mathcal{F})(V)$.
The construction is clearly compatible with relations
(\ref{item-sum}) and (\ref{item-sub})
and compatible with restriction mappings and hence we obtain the map $c$.

\medskip\noindent
Observe that in the discussion above our section $c(s) = \sum (Z_i, s_i)$ of
$f'_!\mathcal{F}$ over $V'$ restricts to zero on
$V' \setminus \Im(\coprod Z_i \to V')$. Since $\Im(\coprod Z_i \to V')$
is proper over $V$ (for example by Morphisms, Lemma
\ref{morphisms-lemma-scheme-theoretic-image-is-proper})
we conclude that $c(s)$ defines a section of
$g_!f'_!\mathcal{F} \subset g_*f'_!\mathcal{F}$ over $V$.
Since every local section of $f_!\mathcal{F}$ locally comes from a
local section of $f_{p!}\mathcal{F}$ we conclude that the image
of $c^\#$ is contained in $g_!f'_!\mathcal{F}$.
Thus we obtain an induced map $c' : f_!\mathcal{F} \to g_!f'_!\mathcal{F}$
factoring $c^\#$ as predicted in the first paragraph of the proof.

\medskip\noindent
Proof of (a). Let $Y'_1 \subset Y'$ be an open subscheme
and set $X_1 = (f')^{-1}(W')$. We obtain a diagram
$$
\xymatrix{
X_1 \ar[d]_{f'_1} \ar[r]_a \ar@/_2em/[dd]_{f_1} &
X \ar[d]^{f'} \ar@/^2em/[dd]^f \\
Y'_1 \ar[d]_{g_1} \ar[r]_{b'} &
Y' \ar[d]^g \\
Y \ar@{=}[r] &
Y
}
$$
where the horizontal arrows are open immersions. Then our claim is that
the diagram
$$
\xymatrix{
f_{1, !}\mathcal{F}|_{X_1} \ar[r]_{c'_1} \ar[dd] &
g_{1, !}f'_{1, !}\mathcal{F}|_{X_1} \ar@{=}[d] \\
& g_{1, !}(f'_!\mathcal{F})|_{Y'_1} \ar[d] \\
f_!\mathcal{F} \ar[r]^{c'} &
g_!f'_!\mathcal{F} \ar[r] & g_*f'_!\mathcal{F}
}
$$
commutes where the left vertical arrow is
Remark \ref{remark-covariance-lqf-f-shriek} and
the right vertical arrow is Remark \ref{remark-covariance-f-shriek-separated}.
The equality sign in the diagram comes about because $f'_1$
is the restriction of $f'$ to $Y'_1$ and our construction
of $f'_!$ is local on the base.
Finally, to prove the commutativity we choose an object $V$ of
$Y_\etale$ and a formal sum $s_1 = \sum (Z_{1, i}, s_{1, i})$ as in
(\ref{equation-formal-sum}) defining an element of
$f_{1, p!}\mathcal{F}|_{X_1}(V)$. Recall this means
$Z_{1, i} \subset X_1 \times_Y V$ is locally closed finite over $V$
and $s_{1, i} \in H_{Z_{1, i}}(\mathcal{F})$.
Then we chase this section
across the maps involved, but we only need to show we
end up with the same element of
$g_*f'_!\mathcal{F}(V) = f'_!\mathcal{F}(Y' \times_Y V)$.
Going around both sides of the diagram the reader immediately
sees we end up with the element $\sum (Z_{1, i}, s_{1, i})$
where now $Z_{1, i}$ is viewed as a locally closed subscheme
of $X \times_{Y'} (Y' \times_Y V) = X \times_Y V$ finite over
$Y' \times_Y V$.

\medskip\noindent
Proof of (b). Let $b : Y_1 \to Y$ be a morphism of schemes. Let us form the
commutative diagram
$$
\xymatrix{
X_1 \ar[d]_{f'_1} \ar[r]_a \ar@/_2em/[dd]_{f_1} &
X \ar[d]^{f'} \ar@/^2em/[dd]^f \\
Y'_1 \ar[d]_{g_1} \ar[r]_{b'} &
Y' \ar[d]^g \\
Y_1 \ar[r]^b &
Y
}
$$
with cartesian squares. We claim that our construction is compatible
with the base change maps of Lemmas \ref{lemma-lqf-base-change-f-shriek}
and \ref{lemma-base-change-f-shriek-separated}, i.e.,
that the top rectangle of the diagram
$$
\xymatrix{
b^{-1}f_!\mathcal{F} \ar[rr] \ar[d]_{b^{-1}c'} & &
f_{1, !}a^{-1}\mathcal{F} \ar[d]^{c_1'} \\
b^{-1}g_!f'_!\mathcal{F} \ar[r] \ar[d] &
g_{1, !}(b')^{-1}f'_!\mathcal{F} \ar[r] \ar[d] &
g_{1, !}f'_{1, !}a^{-1}\mathcal{F} \ar[d] \\
b^{-1}g_*f'_!\mathcal{F} \ar[r] &
g_{1, *}(b')^{-1}f'_!\mathcal{F} \ar[r] &
g_{1, *}f'_{1, !}a^{-1}\mathcal{F}
}
$$
commutes. The verification of this is completely routine and we
urge the reader to skip it. Since the arrows going from the middle
row down to the bottom row are injective, it suffices to show that
the outer diagram commutes.
To show this it suffices to take a local section of
$b^{-1}f_!\mathcal{F}$ and show we end up with the same local
section of $g_{1, *}f'_{1, !}a^{-1}\mathcal{F}$
going around either way. However, in fact it suffices to check
this for local sections which are of the pullback by $b$ of
a section $s = \sum (Z_i, s_i)$ of $f_{p!}\mathcal{F}(V)$
as above (since such pullbacks generate the abelian sheaf
$b^{-1}f_!\mathcal{F}$). Denote $V_1$, $V'_1$, and $Z_{1, i}$
the base change of $V$, $V' = Y' \times_Y V$, $Z_i$ by $Y_1 \to Y$.
Recall that $Z_i$ is a locally closed subscheme of $X_V = X_{V'}$
and hence $Z_{1, i}$ is a locally closed subscheme
of $(X_1)_{V_1} = (X_1)_{V'_1}$. Then $b^{-1}c'$ sends the pullback
of $s$ to the pullback of the local section $c(s) \sum (Z_i, s_i)$ viewed
as an element of $f'_{p!}\mathcal{F}(V') = (g_*f'_{p!}\mathcal{F})(V)$.
The composition of the bottom two base change maps
simply maps this to $\sum (Z_{i, 1}, s_{1, i})$ viewed as an
element of $f'_{1, p!}a^{-1}\mathcal{F}(V'_1) =
g_{1, *}f'_{1, p!}a^{-1}\mathcal{F}(V_1)$.
On the other hand, the base change map at the top of the diagram
sends the pullback of $s$ to $\sum (Z_{1, i}, s_{1, i})$ viewed
as an element of $f_{1, !}a^{-1}\mathcal{F}(V_1)$.
Then finally $c'_1$ by its very construction does indeed
map this to $\sum (Z_{i, 1}, s_{1, i})$ viewed as an
element of $f'_{1, p!}a^{-1}\mathcal{F}(V'_1) =
g_{1, *}f'_{1, p!}a^{-1}\mathcal{F}(V_1)$ and the commutativity
has been verified.

\medskip\noindent
Proof of (c). This follows from comparing the definitions
for both maps; we omit the details.

\medskip\noindent
To finish the proof it suffices to show that the pullback of
$c'$ via any geometric point $\overline{y} : \Spec(k) \to Y$
is an isomorphism. Namely, pulling back by $\overline{y}$
is the same thing as taking stalks and $\overline{y}$
(\'Etale Cohomology, Remark \ref{etale-cohomology-remark-stalk-pullback})
and hence we can invoke
\'Etale Cohomology, Theorem \ref{etale-cohomology-theorem-exactness-stalks}.
By the compatibility (b) just shown, we 
conclude that we may assume $Y$ is the spectrum of $k$
and we have to show that $c'$ is an isomorphism.
To do this it suffices to show that the induced map
$$
\bigoplus\nolimits_{x \in X} \mathcal{F}_x = H^0(Y, f_!\mathcal{F})
\longrightarrow
H^0(Y, g_!f'_!\mathcal{F}) = H^0_c(Y', f'_!\mathcal{F})
$$
is an isomorphism. The equalities hold by
Lemmas \ref{lemma-lqf-f-shriek-stalk} and
\ref{lemma-stalk-f-shriek-separated}.
Recall that $X$ is a disjoint union of
spectra of Artinian local rings with residue field $k$, see
Varieties, Lemma \ref{varieties-lemma-algebraic-scheme-dim-0}.
Since the left and right hand side commute with direct
sums (details omitted) we may assume that $\mathcal{F}$ is a skyscraper
sheaf $x_*A$ supported at some $x \in X$.
Then $f'_!\mathcal{F}$ is the skyscraper sheaf at the
image $y'$ of $x$ in $Y$ by Lemma \ref{lemma-lqf-f-shriek-stalk}.
In this case it is obvious that our construction
produces the identity map $A \to H^0_c(Y', y'_*A) = A$
as desired.
\end{proof}
